\documentclass[a4paper,12pt]{article}
\usepackage[utf8]{inputenc}
\usepackage[english]{babel}
\usepackage{geometry}
\usepackage{fancyhdr}
\usepackage{amsmath,amsfonts,amsthm,amssymb}
\usepackage{graphicx} 
\usepackage{hyperref}
\usepackage{amssymb,graphicx,float,epstopdf,amsmath,amsfonts,amsthm,cases,cite,epsfig,subfigure}
\usepackage[symbol*]{footmisc}

%%%%%%%%%%%%%%%%%%%%%%%%%%%%%%%%%%%%%%%%%%%%%%%%%%%%%%%%%%%%%%%%%%%%%%%%
\topmargin -30pt \advance \topmargin by -\headheight \advance
\topmargin by -\headsep \textheight 9.6in \oddsidemargin -25pt
\evensidemargin \oddsidemargin \marginparwidth 0.5in \textwidth
7.3in

%%%%%%%%%%%%%%%%%%%%%%%%%%%%%%%%%%%%%%%%%%%%%%%%%%%%%%%%%%%%%%%%%%%%%%%%%%%%%%%%
\def\Tr{{\textrm{Tr}}}
\def\str{{\textrm{str}\,}}
\def\tr{{\textrm{tr}}}
\def\diag{{\textrm{diag}}}
\def \eps {\epsilon}

%%%%%%%%%%%%%%%%%%%%%%%%%%%%%%%%%%%%%%%%%%%%%%%%%%%%%%%%%%%%%%%%%%%%%%%%%%%%%%%%

\begin{document}
\section{Boson Star equations}
We take the following scalar field theory
\begin{equation}
S = \int\left[\frac{R + 2\Lambda}{2\kappa} - \frac{1}{4}F_{\mu\nu} F^{\mu\nu} - |D_\mu \phi|^2 - V(|\phi|^2) \right]\sqrt{-g} d^d x,
\end{equation}
where $D_\mu\phi = \partial_\mu\phi - i e A_\mu\phi$.
The equations of motion corresponding to the Lagrangian are
\begin{subequations}
\begin{align}
	&R_{\mu\nu} - \frac{1}{2}R g_{\mu\nu} - \Lambda g_{\mu\nu} = \kappa T_{\mu\nu} = \kappa (T_{\mu\nu}^\phi + T_{\mu\nu}^{em}),\\
	&T_{\mu\nu}^\phi = (D_\mu \phi)^* D_\nu \phi + (D_\nu \phi)^* D_\mu \phi - g_{\mu\nu}(|D_\alpha \phi|^2 + V(|\phi|^2)),\\
	&T_{\mu\nu}^{em} = F_{\mu\alpha} F^{\;\alpha}_\nu - \frac{1}{4} g_{\mu\nu }F_{\alpha\beta}F^{\alpha\beta},\\
	&D^2 \phi - \frac{dV}{d |\phi|^2} \phi = 0,\qquad D^2\phi = g^{\mu\nu} D_\mu D_\nu \phi,\\
	&\triangledown_\mu F^{\mu\nu} = j^\nu = i e (\phi^* D^\nu \phi - (D^\nu\phi)^* \phi).
\end{align}
\end{subequations}
We adopt the following metric ansatz 
\begin{equation}
	\label{eq:AdSAnsatz}
	ds^2 = \frac{l^2}{\cos^2 x}\left(- A(x) e^{-2\delta(x)} dt^2 + A(x)^{-1} dx^2 + \sin^2 x d\Omega^2_{d-1}\right),
\end{equation}
where $l^2 = \frac{d(d-1)}{2 \Lambda}$, the following scalar field ansatz
\begin{equation}
	\phi = \phi_0(x) e^{i\omega t},
\end{equation}
and the following vector potential ansatz
\begin{equation}
	A_\mu = \{a(x),0,\dots\}
\end{equation}
The scalar field equation then becomes
\begin{equation}
	\phi_0'' + \left(\frac{A'}{A} - \delta' + \frac{(d-1)}{\sin x \cos x}\right)\phi_0' + \frac{1}{A}\left(\frac{1}{A}e^{2\delta}(\omega - e a)^2 - \frac{l^2}{\cos^2 x } U_0\right) \phi_0 = 0,\qquad U_0 = \frac{d V}{d |\phi|^2}.
\end{equation}
The gauge field equation becomes
\begin{equation}
	a''(x) + \left(\delta' + \frac{2\cos^2 x + d - 3}{\sin x \cos x}\right) a' + e\frac{2 l^2}{A\cos^2 x} (\omega - e a) \phi_0^2 = 0.
\end{equation}
The stress-energy tensor becomes
\begin{subequations}
\begin{align}
	&\rho = -T^t_t = \frac{\cos^2 x}{l^2}\left(\frac{1}{A}e^{2\delta}(\omega - e a)^2 \phi_0^2 + A(\phi_0')^2\right) + V + \frac{\cos^4 x}{2l^4} e^{2\delta} (a')^2,\\
	&p_r = T^x_x = \frac{\cos^2 x}{l^2}\left(\frac{1}{A}e^{2\delta}(\omega - e a)^2 \phi_0^2 + A(\phi_0')^2\right) - V - \frac{\cos^4 x}{2l^4} e^{2\delta} (a')^2,\\
	&p_\theta = T^i_i = \frac{\cos^2 x}{l^2}\left(\frac{1}{A}e^{2\delta}(\omega - e a)^2 \phi_0^2 - A(\phi_0')^2\right) - V + \frac{\cos^4 x}{2l^4} e^{2\delta} (a')^2,\qquad i \notin \{t, x\}.\\
\end{align}
\end{subequations}

\subsection{Metric decomposition}
Getting the metric part in arbitrary number of dimensions is a harder problem. However, the following trick can be employed. If the metric can be represented in the following form
\begin{equation}
g_{\mu\nu} =\begin{pmatrix}
	f_{ab}(x) & 0 \\
	0 & h_{ij}(y)
\end{pmatrix},
\end{equation}
where $f_{ab}$ is a $d_1$ dimensional metric which is a function of the first $d_1$ coordinates and $h_{ij}$ is $d_2$ dimensional metric which is a function of last $d_2$ coordinates, then the Ricci tensor for this metric trivially becomes
\begin{equation}
	R_{\mu\nu} =\begin{pmatrix}
		R^f_{ab} & 0 \\
		0 & R^h_{ij}
	\end{pmatrix},
\end{equation}
where $R^f_{ab}$ is the Ricci tensor built purely on $f$ metric and $R^h$ is the Ricci tensor built purely on $h$ metric. By applying two successive conformal transformations we can get that the Ricci tensor for the metric 
\begin{equation}
	g_{\mu\nu} =\begin{pmatrix}
		f_{ab}(x) & 0 \\
		0 & e^{2\omega(x)} h_{ij}(y)
	\end{pmatrix},
\end{equation}
where $\omega(x)$ is a function of first $d_1$ coordinates only, can be expressed in a diagonal form
\begin{equation}
R_{\mu\nu} =\begin{pmatrix}
	R^1_{ab} & 0 \\
	0 & R^2_{ij}
\end{pmatrix},
\end{equation}
where
\begin{subequations}
\begin{align}
	&R^1_{ab} = R^f_{ab} - d_2(\triangledown_a \triangledown_b\omega + \partial_a\omega\partial_b\omega),\\
	&R^2_{ij} = R^h_{ij} - e^{2\omega} (\square\omega + d_2(\partial \omega)^2) h_{ij},
\end{align}
\end{subequations}
where all derivatives and all partial derivatives are taken with respect to first $d_1$ coordinates only.

We know that a $d-1$ dimensional sphere has the Ricci tensor
\begin{equation}
R^{sphere}_{ij} = (d - 2) h_{ij},
\end{equation}
therefore an arbitrary spherically symmetric spacetime with a metric
\begin{equation}
	ds^2 = f_{ab} dx^a dx^b + e^{2\omega} d\Omega^2_{d-1},
\end{equation}
has the Ricci tensor
\begin{subequations}
	\begin{align}
		&R^1_{ab} = R^f_{ab} - (d-1)(\triangledown_a \triangledown_b\omega + \partial_a\omega\partial_b\omega),\\
		&R^2_{ij} = ((d-2) - e^{2\omega} (\square\omega + (d-1)(\partial \omega)^2)) h_{ij},
	\end{align}
\end{subequations}
the Einstein tensor is
\begin{subequations}
\label{eq:GEinsteinDecomp}
\begin{align}
	&G^1_{ab} = G^f_{ab} - (d-1)(\triangledown_a\triangledown_b \omega + \partial_a\omega\partial_b \omega) +\\\nonumber
	&\qquad f_{ab} \left((d-1) \square\omega +\frac{1}{2}d(d-1)(\partial\omega)^2 - \frac{1}{2}(d-1)(d-2)e^{-2\omega}\right),\\
	&G^2_{ij} = h_{ij}\left[e^{2\omega}\left((d-2)\square\omega + \frac{1}{2}(d-1)(d-2)(\partial\omega)^2 - \frac{1}{2}R^f\right) - \frac{1}{2}(d-2)(d-3)\right],
\end{align}
\end{subequations}
where $R^f = f^{ab} R_{ab}^f$ is $f$ metric Ricci scalar.

\subsection{Gravitational equations}
Substituting the metric \eqref{eq:AdSAnsatz} into the equations \eqref{eq:GEinsteinDecomp} we obtain the following system of equations that fully describe the AdS Boson Star
\begin{subequations}
\label{eq:StarEOM}
\begin{align}
	&A'(x) = - \frac{d - 2 \cos^2 x}{\sin x \cos x}\left(A - 1\right) - \frac{2 \kappa \sin x \cos x}{d - 1} \left(\frac{e^{2\delta}}{A}(\omega - e a)^2 \phi_0^2 + A(\phi_0')^2 + \frac{\cos^2 x}{2l^2} e^{2\delta} (a')^2\right)\\\nonumber
	&\qquad - \frac{2 \kappa l^2 \tan x}{d - 1} V,\\
	&\delta'(x) = -\frac{2\kappa \sin x \cos x}{d - 1} \left(\frac{e^{2\delta}}{A^2}(\omega - e a)^2 \phi_0^2 + (\phi_0')^2\right),\\
	\label{eq:ScalarEOM}
	&\phi_0'' + \left(\frac{A'}{A} - \delta' + \frac{(d-1)}{\sin x \cos x}\right)\phi_0' + \frac{1}{A}\left(\frac{1}{A}e^{2\delta}(\omega - e a)^2 - \frac{l^2}{\cos^2 x } U_0\right) \phi_0 = 0,\qquad U_0 = \frac{d V}{d |\phi|^2},\\
	&a''(x) + \left(\delta' + \frac{2\cos^2 x + d - 3}{\sin x \cos x}\right) a' + e\frac{2 l^2}{A\cos^2 x} (\omega - e a) \phi_0^2 = 0.
\end{align}
\end{subequations}

\subsection{Point charge solution}
If we assume that no star is present $\phi_0(x) = 0$, the equations \eqref{eq:StarEOM} admit the following solution $\delta(x) = 0$ and
\begin{subequations}
\label{eq:PointCharge}
\begin{align}
	&a(x) = \frac{Q l}{(d-2)}\frac{1}{\tan^{d-2} x},\\
	&A(x) = 1 - \frac{2 M}{l^{d-2}} \frac{\cos^2 x}{\tan^{d-2} x} + \frac{\kappa Q^2}{(d-1)(d-2)} \frac{\cos^2 x}{\tan^{2(d-2)} x}.
\end{align}
\end{subequations}
This solution corresponds to a $(d+1)$-dimensional Reissner-Nordstr\"{o}m black hole with mass $M$ and charge $Q$ in AdS spacetime. By taking $r = l \tan x$ and $\tau = l t$ we can transform this solution to ordinary Schwarzschild coordinates
\begin{equation}
\begin{split}
	&ds^2 = - F(r) d\tau^2 + \frac{dr^2}{F(r)} + r^2 d \Omega^2_{d-1},\\
	&F(r) = 1 - \frac{2M}{r^{d-2}} + \frac{\kappa Q^2}{(d-1)(d-2)}\frac{1}{r^{2(d-2)}} + \frac{r^2}{l^2}.
\end{split}
\end{equation}

\subsubsection{Case $d = 2$}
Case $d = 2$ deserves a special treatment. In this case both the electric potential and the metric become logarithmic in $x$.
\begin{subequations}
	\begin{align}
		&a(x) = Q\log(\tan x),\\
		&A(x) = 1 - 2 M \cos^2 x + \frac{\kappa Q^2}{l^2} (\cos^2 x) \log(\tan x).
	\end{align}
\end{subequations}

\section{Regular solutions for scalar field equation}
In the absence of gravity and electromagnetic field the scalar field equation \eqref{eq:ScalarEOM} becomes
\begin{equation}
	\phi_0'' + \frac{(d-1)}{\sin x \cos x}\phi_0' + \left(\omega^2 - \frac{l^2}{\cos^2 x } U_0\right) \phi_0 = 0,\qquad U_0 = \frac{d V}{d |\phi|^2}.
\end{equation}
It is impossible to solve this equation in general, but it is possible to solve it in the linear case when $V = m^2 |\phi|^2$, the equation then becomes
\begin{equation}
	\phi_0'' + \frac{(d-1)}{\sin x \cos x}\phi_0' + \left(\omega^2 - \frac{l^2 m^2}{\cos^2 x }\right) \phi_0 = 0.
\end{equation}
The general solution of this equation is
\begin{equation}
\begin{split}
	\phi_0 = c_1 (\cos x)^{\alpha_+} F\left(\alpha_+ - \frac{\omega}{2}, \alpha_+ + \frac{\omega}{2}; 2\alpha_+ -\frac{d}{2}+1; \cos^2 x\right) + \\
	c_2 (\cos x)^{\alpha_-} F\left(\alpha_- - \frac{\omega}{2}, \alpha_- + \frac{\omega}{2}; 2\alpha_- -\frac{d}{2}+1; \cos^2 x\right),
\end{split}
\end{equation}
where $\alpha_\pm = \frac{1}{4}(d \pm \sqrt{d^2 + 4 l^2 m^2})$ and $F(a,b;c;x)$ is the hypergeometric function. The second branch is singular at the boundary $x = \pi/2$ so we are going to drop it. We can use the formula
\begin{equation}
\begin{split}
	F(a,b;c;x) =& \frac{\Gamma(c)\Gamma(c-a-b)}{\Gamma(c-a)\Gamma(c-b)}F(a,b;a+b+1-c; 1-x) +\\ &\frac{\Gamma(c)\Gamma(a+b-c)}{\Gamma(a)\Gamma(b)}(1-x)^{c-a-b}F(c-a,c-b;c+1-a-b,1-x),
\end{split}
\end{equation}
find conditions to make the solution regular at the origin. Since $c - a - b = 1 - d/2$, the regularity can only be achieved when $a = - j$ and $j$ is zero or a positive integer. In such case $1/\Gamma(a)$ becomes zero. Therefore if we take
\begin{equation}
	\omega = 2\alpha_+ + 2j,
\end{equation}
we obtain a solution regular everywhere
\begin{equation}
	\phi_0 = c_1 (\cos x)^{2\alpha_+} F\left(-j, 2\alpha_+ + j; 2\alpha_+ -\frac{d}{2}+1; \cos^2 x\right),
\end{equation}
the Hypergeometric function is a polynom in $\cos^2 x$. It is more convenient to rewrite this expression in terms of Jacobi polynomials
\begin{equation}
	\phi_j = C (\cos x)^{(d + \sqrt{d^2 + 4 l^2 m^2})/2} P^{d/2 - 1, \sqrt{d^2 + 4 l^2 m^2}/2}_j(\cos 2x),
\end{equation}
the corresponding frequency explicitly is
\begin{equation}
	\omega_j = \frac{1}{2}(d + \sqrt{d^2 + 4 l^2 m^2}) + 2j.
\end{equation}
The fundamental mode corresponds to $j = 0$
\begin{equation}
	\phi_0 = C (\cos x)^{\frac{1}{2}(d + \sqrt{d^2 + 4 l^2 m^2})}
\end{equation}

\section{Asymptotic expansion}
\subsection{Boundary case $x = \pi/2$}
We introduce a new coordinate $\rho = \pi/2 - x$. If we expand the solution in terms of $\rho$ around $\rho = 0$ we get
\begin{subequations}
\begin{align}
	&A(\rho) = 1 + A^{(0)}\left(1 + \frac{d-3}{3}\rho^2\right) \rho^d + (a^{(1)})^2\frac{\kappa (d-2) e^{2\delta_0}}{l^2 (d-1)}\left(1 + \frac{2d-5}{3}\rho^2\right) \rho^{2d-2} + \\\nonumber
	&\qquad\qquad (\phi^{(0)})^2 \frac{2\kappa n^2}{(2n - d)(d-1)} \rho^{2n} + O(\rho^{d+4}) + O(\rho^{2d+2}) + O(\rho^{2n+2}),\\
	&\delta(\rho) = \delta_0 + \frac{\kappa n(\phi^{(0)})^2}{d-1} \rho^{2n} + O(\rho^{2n+2}),\\
	&\phi_0(\rho) = \phi^{(0)} \rho^{n} \left(1 + \frac{n(n + d - 2) - 3e^{2\delta_0}(\omega - e a^{(0)})^2}{6(2(n+1) - d)} \rho^2 - \frac{A^{(0)}\rho^d}{2}\right) + O(\rho^{n+4}),\\
	&a(\rho) = a^{(0)} + a^{(1)}\left(1 + \frac{d-2}{3}\rho^2\right)\rho^{d-2} + O(\rho^{d+2}).
\end{align}
\end{subequations}
where we label $V'(0) = m^2$ and
\begin{equation}
	n = \frac{d}{2} + \frac{1}{2}\sqrt{d^2 + 4 l^2 m^2}
\end{equation}
We compare the expressions with asymptotic expressions for point charge and see that $A^{(0)}$ corresponds to the star mass and $a^{(1)}$ corresponds to the star charge
\begin{equation}
	A^{(0)} = - \frac{2 M}{l^{d-2}},\qquad a^{(1)} =\frac{Q l}{d-2}.
\end{equation}

We now develop a method to reconstruct asymptotic coefficients from function values at a certain point $\rho$. We introduce the following notations
\begin{equation}
	A_0 = A^{(0)}\rho^d,\qquad a_0 = a^{(0)},\qquad a_1 = a^{(1)}\rho^{d-2},\qquad \phi_0 = \phi^{(0)}\rho^n.
\end{equation}
We begin by reconstructing $a_0$ and $a_1$ as
\begin{subequations}
\begin{align}
	&a_1 = \frac{\rho a'(\rho)}{(d-2)\left(1 + \frac{d}{3}\rho^2\right)},\\
	&a_0 = a(\rho) - a_1\left(1 + \frac{d-2}{3}\rho^2\right).
\end{align}
\end{subequations}
Within the assumed order of approximation we can take
\begin{equation}
	e^{2\delta_0} = e^{2\delta(\rho)}\left(1 - \frac{2\kappa n}{d-1} \phi_0^2\right).
\end{equation}
We can obtain $A_0$ from the expression
\begin{equation}
\label{eq:A_asym_expr}
	A(\rho) = 1 + A_0\left(1 + \frac{d-3}{3}\rho^2\right) + e^{2\delta_0} b_4 a_1^2 \left(1 + \frac{2d-5}{3}\rho^2\right) \rho^{2} + b_5\phi_0^2,
\end{equation}
where
\begin{equation}
	b_4 = \frac{\kappa (d-2)}{l^2 (d-1)},\qquad b_5 = \frac{2\kappa n^2}{(2n - d)(d-1)}.
\end{equation}

Substituting this expression into $\phi_0(\rho)$ asymptotic and taking the leading asymptotic to get $A^0$ we get a cubic equation for $\phi_0$
\begin{equation}
	\phi_0(\rho) = \phi_0 \left(1 + b_1\rho^2 - e^{2\delta_0} b_2\rho^2 - \frac{1}{2} A_0\right)
\end{equation}
where
\begin{equation}
	b_1 = \frac{n(n + d - 2)}{6(2(n+1) - d)},\qquad b_2 = \frac{(\omega - e a_0)^2}{2(2(n+1) - d)},\qquad b_3 = \frac{\kappa n}{d-1}.
\end{equation}
we get a cubic equation for $\phi_0$
\begin{equation}
\begin{split}
	&\left(2 b_2 b_3 e^{2\delta(\rho)}\rho^2 + \frac{b_5}{2(1 + (d-3)\rho^2/3)}\right)\phi_0^3 + \\
	&\qquad\left(1 + (b_1 - b_2 e^{2\delta(\rho)})\rho^2 - \frac{A(\rho) - 1 - e^{2\delta(\rho)} b_4 a_1^2 (1 + (2d - 5)\rho^2/3)\rho^2}{2(1 + (d-3)\rho^2/3)}\right)\phi_0 = \phi_0(\rho),
\end{split}
\end{equation}
This cubic equation can be solved for $\phi_0$ after which $\delta_0$ is obtained as
\begin{equation}
	\delta_0 = \delta(\rho) - b_3 \phi_0^2,
\end{equation}
and $A_0$ can be found from the expression \eqref{eq:A_asym_expr}.
\end{document}